\section*{摘要}
针对含光伏发电与储能装置的社区微网购电调度问题，考虑负荷、光伏出力及电价在不同决策时点下的不确定性，本文以实际结算成本最小为目标，以 $10\,\mathrm{min}$ 为调度粒度，构建日前购电决策、日内滚动执行与跨日状态传递相结合的协同优化框架，并基于 2025 年数据开展因果回测。

针对问题一，在预测信息已知条件下，构建日前购电与储能协同调度的确定性线性规划模型。模型综合考虑供需平衡、储能充放电效率、SOC 边界、充放电功率及首末储能状态一致性等约束，采用 HiGHS 求解日前最优购电与储能计划。结果表明，购电费用为 $33801.50$ 元，较无储能方案节省 $14250.55$ 元，降幅为 $29.66\%$；储能能够通过峰谷套利和提高光伏消纳水平降低购电成本。

针对问题二，在预测偏差引起的计划购电与高价紧急购电权衡条件下，构建基于 K-medoids 场景压缩、风险储备和滚动模型预测控制的日前与日内协同调度模型。日初利用历史信息生成计划购电额度，日内依据已观测误差滚动修正剩余时段的储能决策与紧急购电量，并以 8 个代表场景刻画负荷与光伏预测误差。对 2025 年 2 月至 12 月共 334 日进行因果回测，实际购电总成本为 $15166538.46$ 元，其中计划购电成本为 $11883687.24$ 元，紧急购电成本为 $3282851.22$ 元。该模型实现了计划购电成本与供电缺口风险之间的协调。

针对问题三，在多时点光伏预测下既有购电承诺动态调整条件下，构建包含日初承诺、信息价值评估及剩余时段重优化的滚动购电调整模型。在 6:00、12:00 和 18:00 获取新预测后，模型锁定已执行时段，仅当更新预测带来的预期收益超过调整费用时，才修改未执行时段的购电计划。全年连续运行结果表明，价值触发调整策略的总成本为 $16373508.75$ 元，较不调整策略降低 $736059.50$ 元，降幅为 $4.30\%$，表明该策略在本回测期内能够降低实际结算成本。

针对问题四，在实时波动电价条件下，基于由负荷、光伏和电价构成的三元联合残差场景，构建日前决策与日内滚动执行相结合的优化模型。日前阶段利用代表场景制定购电额度，日内阶段依据实际价格和已观测信息滚动更新剩余时段决策。结果表明，日前额度策略在 2 月至 12 月的实际结算成本为 $15257873.75$ 元；多时点价值触发调整策略的同期成本为 $15346418.14$ 元。由于两类策略的购电承诺、调整费用及紧急购电结算机制不同，比较时应在统一信息边界与结算区间下进行。

最后，对各模型进行核验，并评估模型优缺点。

\par\noindent 关键词：社区微网；储能调度；线性规划；模型预测控制；因果回测；购电调整
